\documentclass[a4paper,11pt]{article}
%% Compiles both from this folder (pdflatex tp1.tex) and from the root of the project
%% (Overleaf: Menu > Main document = tp1_protocol/tp1.tex).
\makeatletter
\IfFileExists{tp1_protocol/tp1.tex}{\gdef\tpdir{tp1_protocol/}\gdef\tpcommon{common/}}{\gdef\tpdir{}\gdef\tpcommon{../common/}}
\edef\input@path{{\tpdir}{\tpcommon}}
\makeatother
\input{tpsetup}
\renewcommand{\polytitle}{Statistical Supervised Machine Learning -- TP 1}
\graphicspath{{figures/}{tp1_protocol/figures/}{../common/logos/}{common/logos/}}

\begin{document}
\sloppy
\sheethead{TP 1 -- An experimental protocol for supervised learning}{%
This first practical session sets up the experimental protocol that we will use in all the following ones. On two
simple families of models, Gaussian kernel SVMs and decision trees, we distinguish the empirical risk from the true
risk, we see why and when the data must be normalized, we search hyperparameters with a validation set and with
cross-validation, and we evaluate the final model once, on a test set kept aside. The last two parts open the
programme of the course: combining many trees into a vote (ensemble methods, TP~2 and~3), and turning a test error into
a \emph{guarantee} on the true risk (generalization bounds and PAC-Bayes, TP~4 and~5). The session lasts
\textbf{three hours}; Part~\ref{p1:far} contains the difficult questions and can be left for home.}

\begin{sessionbox}[About this session -- TP 1: An experimental protocol]
\textbf{Goal.} By the end of the session you should be able to explain the difference between empirical and true
risk, decide when to normalize the features, tune hyperparameters by cross-validation without touching the test set,
report a test error with its uncertainty, and say what ensemble methods and generalization bounds will bring.\par
\textbf{To reread before the session.} Sections~5 and~6.1 (risk and empirical risk minimization), Section~7.1 (how to
train and tune a model), Section~8.3 (SVM and kernels), Section~8.7 (decision trees) and Sections~11.3--11.5
(preprocessing, splitting, tuning) of the lecture notes.\par
\textbf{Timing (3 hours).}\quad
\begin{tabular}[t]{ll}Setting up and Part~1 & 40 min\\
Part~2 & 20 min\\
Part~3 & 45 min\\
Part~4 & 25 min\\
Part~5 & 20 min\\
Part~6 & 30 min\\
Part~7 \hard & at home\end{tabular}\par
\textbf{To hand in.} Your notebook, with the figures of Parts~1 and~3 and the tables of Parts~2 to~6.\par
\textbf{Tools.} \code{numpy}, \code{matplotlib} and \code{scikit-learn} (\code{SVC}, \code{DecisionTreeClassifier},
\code{StandardScaler}, \code{make\_pipeline}, \code{StratifiedKFold}, \code{cross\_val\_score}, \code{GridSearchCV}).
\end{sessionbox}

\medskip
\noindent\textbf{Data.} A synthetic problem whose distribution is known (Part~\ref{p1:risk}), and three datasets of the
course loaded with \code{load(name)} of \code{common/tptools.py}, with labels in $\{-1,+1\}$: \code{wdbc} (breast cancer,
$569\times30$), \code{pima} (diabetes, $768\times8$) and \code{australian} (credit, $690\times14$). The file
\code{tp1\_starter.py} contains the data generator and the imports.

\begin{tip}
Fix all the random seeds (\code{random\_state=seed}). Never call \code{fit} on data that contains test examples, and
put every preprocessing step in a \code{Pipeline} so that cross-validation refits it on each training fold. Keep the
test set in a variable that you use only once, at the end.
\end{tip}

%% =====================================================================
\tppart{Empirical risk and true risk (Sections 5 and 6.1)}\label{p1:risk}

The input $X$ is uniform on $[-1,1]^2$, the clean label is $c(x)=+1$ if $x_2>\frac12\sin(\pi x_1)$ and $-1$ otherwise,
and the observed label $Y$ is $c(X)$ flipped with probability $\eta=0.1$. For a classifier $h$, the true risk is
$\RD(h)=\P(h(X)\neq Y)$ and the empirical risk on a sample $S=\{(x_i,y_i)\}_{i=1}^m$ is
$\RS(h)=\frac1m\sum_{i=1}^m\ind{h(x_i)\neq y_i}$.

\begin{enumerate}
\item What is the Bayes classifier of this problem and its risk $R^\star$?\\
      Check it on $50\,000$ examples.
\item Draw a training sample of size $m=200$ and a sample of $50\,000$ examples used as a proxy for the true risk.\\
      Plot the empirical (training) risk and the true risk of a decision tree as functions of its maximal depth
      ($1$ to $15$), and of a Gaussian kernel SVM ($C=10$) as functions of $\gamma\in[10^{-2},10^{3}]$, where the kernel
      is $k(x,x')=e^{-\gamma\|x-x'\|^2}$.\\
      Identify underfitting and overfitting.\\
      Plot the decision boundaries of the
      SVM for $\gamma\in\{0.01,3,100\}$.
\item Fix the tree of depth $5$.\\
      Draw $500$ fresh samples of size $n\in\{50,200,1000,5000\}$, compute its empirical
      risk on each, and compare their mean and standard deviation with $\RD(h)$ and $\sqrt{\RD(h)(1-\RD(h))/n}$.\\
      Why is the empirical risk an unbiased estimate of the true risk here, but not in question~2?
\end{enumerate}

%% =====================================================================
\tppart{Normalization of the data (Section 11.3)}\label{p1:norm}

\begin{enumerate}
\item Look at the standard deviations of the features of \code{wdbc} and \code{australian}.\\
      What do you notice?
\item For $10$ random splits ($70\%/30\%$), compute the test error of an SVM (\code{C=1}, \code{gamma="scale"}) and of a
      tree of depth $4$, with and without standardization (\code{StandardScaler} fitted on the training set).\\
      Explain the
      results.
\item Why must the scaler be fitted on the training set only?\\
      How do you do it inside a cross-validation?
\end{enumerate}

%% =====================================================================
\tppart{Hyperparameter search and cross-validation (Sections 7.1 and 11.5)}\label{p1:cv}

We work on \code{pima}. Keep $25\%$ of the data as test set (stratified, \code{random\_state=0}); the remaining
$75\%$ is the learning set.

\begin{enumerate}
\item On the learning set, compute the $5$-fold cross-validation error of the pipeline (standardization + SVM) for
      $C\in\{0.1,1,10,100,1000\}$ and $\gamma\in\{10^{-4},\dots,1\}$, and display it as a heat map.\\
      Do the same for a
      tree, with \code{max\_depth}$\in\{1,2,3,4,6,8,\text{None}\}$ and \code{min\_samples\_leaf}$\in\{1,5,10,20\}$.\\
      Which hyperparameters do you select?\\
      Retrain on the whole learning set and compute the test error.
\item \textbf{Hold-out or cross-validation?} Repeat the selection of $(C,\gamma)$ $30$ times, with different random
      partitions of the learning set, (a) with a single hold-out validation split ($75\%/25\%$ of the learning set),
      (b) with $5$-fold cross-validation.\\
      For each procedure, count the number of distinct pairs selected, and compute
      the mean and standard deviation of the selected validation score and of the test error of the selected model.\\
      Conclude.
\end{enumerate}

%% =====================================================================
\tppart{The complete protocol, and the optimism of the best score}\label{p1:protocol}

\begin{enumerate}
\item Write a function that, given a dataset and a seed, (i) splits it into a learning set ($75\%$) and a test set,
      (ii) selects the hyperparameters of the SVM pipeline and of the tree by $5$-fold cross-validation on the learning
      set (\code{GridSearchCV}), (iii) retrains them on the whole learning set and (iv) computes their test errors.\\
      Run it on \code{wdbc}, \code{pima} and \code{australian} for $10$ seeds, and compare the best cross-validation
      error with the test error.
\item To understand the difference, replace the labels of \code{pima} by random labels (independent of the features)
      and run the same protocol for the SVM.\\
      What is the true risk of any classifier?\\
      What do you observe?
\end{enumerate}

%% =====================================================================
\tppart{Opening: what if we combine many trees? (Section 10.2)}\label{p1:ens}

\begin{enumerate}
\item On the synthetic problem ($m=300$) and on \code{pima} ($70\%/30\%$), learn $B=100$ fully grown trees, each one on a
      \emph{bootstrap sample} of the training set ($m$ indices drawn with replacement), and predict by the majority vote
      $\sign\big(\sum_{b}h_b(x)\big)$.\\
      Compare the test error of the vote with that of one fully grown tree learned on the
      whole training set, and with the mean test error of the $100$ trees.
\item Why can a vote be much better than its voters?
\end{enumerate}

%% =====================================================================
\tppart{Opening: a first guarantee on the true risk (Sections 1.4 and 6.2)}\label{p1:bounds}

Let $h$ be a classifier \emph{fixed before} drawing a test sample $T$ of $n$ i.i.d.\ examples. Hoeffding's inequality
gives, with probability at least $1-\delta$ over $T$,
\[
\RD(h)\leq\widehat R_T(h)+\sqrt{\frac{\ln(1/\delta)}{2n}} ,
\]
and a tighter version, based on the Chernoff bound, is $\RD(h)\leq\klup\big(\widehat R_T(h),\frac{\ln(1/\delta)}{n}\big)$,
where $\klup(q,\varepsilon)=\sup\{p\in[q,1]:\kl(q\|p)\leq\varepsilon\}$ and $\kl(q\|p)=q\ln\frac qp+(1-q)\ln\frac{1-q}{1-p}$
(\code{kl\_inv\_upper} in \code{common/tptools.py}). We take $\delta=0.05$.

\begin{enumerate}
\item Prove the Hoeffding bound from $\P\big(\RD(h)-\widehat R_T(h)\geq\varepsilon\big)\leq e^{-2n\varepsilon^2}$.
\item Compute both bounds for the SVM selected in Part~\ref{p1:protocol} (seed $0$) on the three datasets.\\
      Interpret them.\\
      How many test examples would be needed for a Hoeffding term of $0.01$?
\item On the synthetic problem, fix the SVM with $\gamma=3$ learned on $200$ examples, draw $1000$ test sets of size
      $n\in\{50,200,1000\}$ and check that both bounds hold in at least $95\%$ of the cases.
\item \textbf{Selection.} Why can the bound not be applied to the best cross-validation score of Part~\ref{p1:protocol}?\\
      Simulate $K$ useless classifiers (true risk $\frac12$) whose errors on a common validation set of $n=100$ examples
      are independent $\mathcal B(n,\frac12)/n$, select the best one, and compute how often the Hoeffding bound holds for
      it, for $K\in\{1,25,1000\}$.\\
      Show that replacing $\delta$ by $\delta/K$ (union bound) repairs it.
\end{enumerate}

%% =====================================================================
\tppart{Going further \hard}\label{p1:far}

\begin{enumerate}
\item Prove that the Bayes classifier $h^\star(x)=\sign(2\eta(x)-1)$ (either label when $\eta(x)=\frac12$), with $\eta(x)=\P(Y=1\mid X=x)$, minimizes the true
      risk, and that $R^\star=\E[\min(\eta(X),1-\eta(X))]$.
\item \textbf{Data leakage.} Draw $m=100$ examples with $d=2000$ Gaussian features and labels drawn uniformly in
      $\{-1,+1\}$, independently of the features.\\
      (a) Select the $20$ features most correlated with the label on the
      whole dataset, then estimate the accuracy of a logistic regression on these features by $5$-fold cross-validation.\\
      (b) Do the same, but perform the selection inside each fold, on the training part only.\\
      Repeat $10$ times and
      compare with the true accuracy.
\end{enumerate}

\end{document}
